\chapter{Agricoles et softs}\label{ch:m3:agriculturals-and-softs}

À midi, heure de l'Est, le 30 juin 2026, le US Department of Agriculture a publié son estimation des
surfaces que les agriculteurs américains avaient semées ce printemps~: le maïs recule de 3~\% sur un an, le soja progresse de 5~\%.
Pour les traders qui suivent le marché, les minutes qui suivent midi étaient les plus importantes de la saison~: le
chiffre leur indiquait la quantité de maïs qui existerait à l'automne, et les contrats à terme le réévaluaient avant que la plupart
des agriculteurs aient lu le communiqué. Les marchés agricoles fonctionnent selon des calendriers que la nature fixe et
que les gouvernements mesurent. Ce chapitre traite de la \omterm{def:m3:agriculturals-and-softs:cropyear}{campagne} et de ses spreads, des rapports qui font bouger le
marché, la météo, les limites quotidiennes qui peuvent bloquer un trader, et le rapport hebdomadaire qui montre
qui détient quoi.

\section{Calendriers de culture et spread entre ancienne et nouvelle récolte}

\begin{definition}[Campagne, ancienne récolte, nouvelle récolte]\label{def:m3:agriculturals-and-softs:cropyear}
La \emph{campagne}\index{campagne} (marketing year) d'une culture est la période de douze mois
au cours de laquelle une récolte est vendue et utilisée~; pour le maïs américain, elle court du 1\up{er} septembre au 31 août.
L'\emph{ancienne récolte}\index{ancienne récolte} désigne le grain issu de la récolte déjà en stock~; la \emph{nouvelle
récolte}\index{nouvelle récolte} désigne le grain de la récolte à venir. Les contrats livrables avant la nouvelle récolte
se cotent sur l'ancienne récolte, ceux livrables après sur la nouvelle.
\end{definition}

\begin{omfigure}
\centering
\begin{tikzpicture}[font=\small,x=0.9cm]
\draw[thick] (0,0) -- (12,0);
\foreach \i/\m in {0/janv.,1/févr.,2/mars,3/avr.,4/mai,5/juin,6/juil.,7/août,8/sept.,9/oct.,10/nov.,11/déc.}{
  \draw (\i,0.1) -- (\i,-0.1); \node[anchor=north,font=\footnotesize] at (\i+0.5,-0.1) {\m};}
\fill[omProp!35] (3.5,0.3) rectangle (5.5,0.8); \node[font=\footnotesize] at (4.5,0.55) {semis};
\fill[omThm!35] (6.4,0.3) rectangle (7.3,0.8); \node[font=\footnotesize,anchor=south west] at (6.4,0.85) {pollinisation};
\fill[omDef!35] (8.5,0.3) rectangle (10.8,0.8); \node[font=\footnotesize] at (9.65,0.55) {récolte};
\draw[{Stealth[length=2mm]}-{Stealth[length=2mm]},thick] (8,-0.9) -- (12,-0.9) node[midway,below,font=\footnotesize] {nouvelle campagne};
\draw[thick,-{Stealth[length=2mm]}] (0,-0.9) -- (7.9,-0.9) node[midway,below,font=\footnotesize] {ancienne récolte vendue sur stock};
\node[font=\footnotesize,anchor=south east] at (5.95,1.1) {30 juin~: surfaces};
\draw[gray,dashed] (5.97,0.2) -- (5.97,1.1);
\end{tikzpicture}

\omcaption{Le maïs américain sur une année, schéma~: semis au printemps, pollinisation en juillet
(les semaines sensibles à la météo), récolte à partir de l'automne~; la \omterm{def:m3:agriculturals-and-softs:cropyear}{campagne} commence le 1\up{er} septembre. Le rapport
sur les surfaces de fin juin est la première donnée ferme de la saison sur la taille de la \omterm{def:m3:agriculturals-and-softs:cropyear}{nouvelle récolte}.}\label{fig:m3:agriculturals-and-softs:calendar}
\end{omfigure}

Le spread entre le dernier contrat à terme sur l'\omterm{def:m3:agriculturals-and-softs:cropyear}{ancienne récolte} et le premier sur la nouvelle (pour le maïs, juillet contre
décembre) est un pari sur la taille de la \omterm{def:m3:agriculturals-and-softs:cropyear}{nouvelle récolte} face aux stocks laissés par l'ancienne. Une \omterm{def:m3:agriculturals-and-softs:cropyear}{ancienne
récolte} faible et une \omterm{def:m3:agriculturals-and-softs:cropyear}{nouvelle récolte} abondante placent juillet bien au-dessus de décembre~; un excédent de vieux grain et une \omterm{def:m3:agriculturals-and-softs:cropyear}{nouvelle récolte} défaillante
inversent la structure.

\section{Rapports publics}

\begin{definition}[Rapport WASDE, ratio stocks/utilisation]\label{def:m3:agriculturals-and-softs:wasde}
Le \emph{rapport WASDE}\index{rapport WASDE} (World Agricultural Supply and Demand Estimates) est le bilan mensuel publié par le
US Department of Agriculture, qui recense l'offre, l'utilisation et les stocks finaux des principales cultures,
aux États-Unis et dans le monde. Le \emph{ratio stocks/utilisation}\index{ratio stocks/utilisation} est le
rapport entre les stocks finaux d'une \omterm{def:m3:agriculturals-and-softs:cropyear}{campagne} et l'utilisation totale de cette même \omterm{def:m3:agriculturals-and-softs:cropyear}{campagne}~: la réserve que le marché
conserve jusqu'à la récolte suivante.
\end{definition}

Le prix d'une culture stockable dépend de sa rareté en fin de \omterm{def:m3:agriculturals-and-softs:cropyear}{campagne}. Quand le
\omterm{def:m3:agriculturals-and-softs:wasde}{ratio stocks/utilisation} est faible, la réserve est mince et le prix doit rationner la consommation~; un ratio élevé
signifie qu'il faut payer le stockage. Les rapports qui révisent le bilan font bouger le prix aussitôt.

\begin{dated}{2026-09}{Calendrier des publications}\label{dat:m3:agriculturals-and-softs:reports}
Le \omterm{def:m3:agriculturals-and-softs:wasde}{rapport WASDE} est publié à 12~h, heure de l'Est, en général entre le 8 et le 12 de chaque mois. Le
rapport Acreage du National Agricultural Statistics Service est publié à midi, heure de l'Est, à la fin
de juin~; le 30 juin 2026, il fait ressortir des surfaces de maïs semées inférieures de 3~\% à celles de 2025 et des surfaces de soja supérieures de 5~\%.
\end{dated}

\begin{example}[Une surprise sur le ratio stocks/utilisation]\label{ex:m3:agriculturals-and-softs:stu}
Un bilan comptant 15.0 milliards de boisseaux d'utilisation et 1.8 milliard de stocks finaux affiche un
\omterm{def:m3:agriculturals-and-softs:wasde}{ratio stocks/utilisation} de 12.0~\%. Un rapport qui réduit la récolte de 300 millions de boisseaux, à utilisation inchangée, ramène
les stocks finaux à 1.5 milliard~: 10.0~\%. Le prix doit monter jusqu'à ce que l'utilisation recule assez pour reconstituer la réserve que
le marché souhaite conserver.
\end{example}

\section{Météo}

Pour une culture en pleine croissance, la météo est le risque dominant, et il se concentre sur quelques semaines~: pour le maïs,
pollinisation en juillet, quand la chaleur et la sécheresse réduisent le plus les rendements. Les traders achètent ce risque avant les
semaines et le revendent ensuite, un schéma appelé la prime météo~; les prévisions des principaux
modèles numériques font bouger les prix pendant la saison aussi sûrement que le font les rapports.

\section{Limites quotidiennes}

Les contrats à terme sur céréales sont soumis à des limites quotidiennes de variation de cours~: un contrat ne peut pas se traiter à plus d'un montant fixé au-dessus ou
en dessous du règlement-livraison de la veille. La limite est une limite de variation de cours du type décrit au livre~1, chapitre~12,
pour les actions.

\begin{definition}[Marché bloqué en limite, limite élargie]\label{def:m3:agriculturals-and-softs:limit}
Un \emph{marché bloqué en limite}\index{marché bloqué en limite} est un marché dont le cours a atteint la limite quotidienne
avec des ordres d'un seul côté~: aucune transaction ne peut se conclure à un cours au-delà de cette limite, si bien que les positions ne peuvent pas être clôturées.
Une \emph{limite élargie}\index{limite élargie} est la limite plus large qui s'applique le lendemain
d'un contrat dont le règlement-livraison s'est fait à sa limite.
\end{definition}

\begin{dated}{2026-09}{Limites sur les céréales}\label{dat:m3:agriculturals-and-softs:limits}
Contrats à terme CBOT sur le maïs~: 5\,000 boisseaux, cotés en cents par boisseau. Leurs limites sont réajustées tous les six mois~:
7~\% du règlement-livraison moyen d'un contrat de référence sur 45 jours de cotation, arrondi
au multiple de 5 cents le plus proche, avec un plancher de 20 cents~; la \omterm{def:m3:agriculturals-and-softs:limit}{limite élargie} vaut 1.5 fois cette valeur, arrondie au multiple supérieur. Pour une moyenne de
4.50~\$, la limite est de 30 cents et la \omterm{def:m3:agriculturals-and-softs:limit}{limite élargie} de 45 cents. Pour la date de transaction du 8 septembre 2026, le CME
Group a fixé la limite sur le maïs à 30 cents (60 cents pour les calendar spreads).
\end{dated}

\begin{proposition}[Marge pour survivre à un blocage]\label{prop:m3:agriculturals-and-softs:lock}
Si une information fait varier de $\Delta$ la fair value d'un contrat au-delà de sa limite $L$, la \omterm{def:m3:agriculturals-and-softs:limit}{limite élargie} étant $L_e$,
un détenteur placé du mauvais côté ne peut pas sortir tant que le cours n'a pas rejoint la fair value~: il perd $L$ le premier
jour et $L_e$ chaque jour de blocage supplémentaire, et il lui faut une marge pour chaque jour bloqué avant de pouvoir négocier.
Avec $L < |\Delta| \le L + L_e$ le blocage dure deux jours.
\end{proposition}

\begin{proof}
Chaque jour bloqué, le règlement-livraison se déplace de la limite du jour et aucune transaction n'est possible à la fair value~;
le blocage prend fin le premier jour où l'écart restant tient dans la limite.
\end{proof}

\begin{omfigure}
\begin{tikzpicture}
\begin{axis}[width=11cm,height=4.4cm,
  xlabel={jour de négociation},ylabel={\$ par boisseau},
  tick label style={font=\small},label style={font=\small},
  xmin=-0.3,xmax=5.3,ymin=3.5,ymax=4.7,grid=major,grid style={gray!25},xtick={0,1,2,3,4,5},
  legend columns=2,legend style={font=\footnotesize,at={(0.5,-0.33)},anchor=north,draw=none,fill=none,/tikz/every even column/.append style={column sep=8pt}},
  table/col sep=comma]
\addplot[omThm,very thick,mark=*,mark size=1.6pt] table[x=day,y=settle]{figdata/markets-3/09-agriculturals-and-softs/limits.csv};
\addplot[omDef,thick,dashed,const plot] table[x=day,y=fair]{figdata/markets-3/09-agriculturals-and-softs/limits.csv};
\node[font=\footnotesize,anchor=west] at (axis cs:1.05,4.23) {limite basse de 30 cents, bloquée};
\node[font=\footnotesize,anchor=west] at (axis cs:2.05,3.83) {limite élargie de 45 cents, bloquée};
\legend{règlement-livraison,fair value après le rapport}
\end{axis}
\end{tikzpicture}

\omcaption{Maïs après un rapport qui ramène sa fair value de 4.50~\$ à 3.70~\$ le boisseau~: deux jours bloqués
à la limite et à la \omterm{def:m3:agriculturals-and-softs:limit}{limite élargie}, puis négociation à la fair value. Illustratif~; limites selon la
règle de l'encadré daté. Données~: le tutoriel du chapitre.}\label{fig:m3:agriculturals-and-softs:limits}
\end{omfigure}

\section{Rapports de positionnement}

\begin{definition}[rapport Commitments of Traders, managed money]\label{def:m3:agriculturals-and-softs:cot}
Le \emph{rapport Commitments of Traders}\index{rapport Commitments of Traders} (COT) est la publication hebdomadaire
de la US Commodity Futures Trading Commission (CFTC) sur les positions longues et courtes agrégées
des grands intervenants de chaque marché à terme, par catégorie, à la date du mardi. Sous sa forme désagrégée,
\emph{gestion pour compte de tiers}\index{gestion pour compte de tiers} (managed money) désigne la catégorie des commodity trading advisors, commodity
pool operators et des fonds qui négocient pour le compte de clients, à la différence des producteurs et négociants,
des swap dealers et des autres reportable traders.
\end{definition}

\begin{dated}{2026-09}{Calendrier du COT}\label{dat:m3:agriculturals-and-softs:cot}
La Commission publie le rapport chaque vendredi à 15:30, heure de l'Est, avec les positions arrêtées au
mardi précédent.
\end{dated}

\begin{omfigure}
\begin{tikzpicture}
\begin{axis}[width=11cm,height=4.8cm,axis y line*=left,
  ylabel={managed money net, \% des positions ouvertes},
  tick label style={font=\small},label style={font=\small},
  xmin=2016,xmax=2026.8,ymin=-30,ymax=30,grid=major,grid style={gray!25},
  xtick={2016,2018,2020,2022,2024,2026},xticklabel style={/pgf/number format/1000 sep={}},
  table/col sep=comma]
\addplot[omDef,very thick] table[x=t,y=mm]{figdata/markets-3/09-agriculturals-and-softs/mm.csv};\label{plot:m3:mm}
\draw[gray] (axis cs:2016,0) -- (axis cs:2026.8,0);
\end{axis}
\begin{axis}[width=11cm,height=4.8cm,axis y line*=right,axis x line=none,
  ylabel={maïs, \$/t},ylabel near ticks,
  tick label style={font=\small},label style={font=\small},
  xmin=2016,xmax=2026.8,ymin=100,ymax=400,
  legend columns=2,legend style={font=\footnotesize,at={(0.5,-0.18)},anchor=north,draw=none,fill=none,/tikz/every even column/.append style={column sep=8pt}},
  table/col sep=comma]
\addlegendimage{/pgfplots/refstyle=plot:m3:mm}\addlegendentry{managed money net (gauche)}
\addplot[omProp,thick,dashed] table[x=t,y=maize]{figdata/markets-3/09-agriculturals-and-softs/mm.csv};\addlegendentry{prix du maïs (droite)}
\end{axis}
\end{tikzpicture}

\omcaption{Maïs CBOT~: position nette du managed money en part des positions ouvertes, moyennes mensuelles des
rapports hebdomadaires, de janvier 2016 à juillet 2026, comparée au prix mensuel du maïs. Les fonds sont acheteurs quand les prix
montent et vendeurs quand ils baissent~: les deux séries ont une corrélation de 0.68. Données~: CFTC, Commitments of Traders désagrégé
(contrats à terme uniquement)~; prix du maïs du FMI via FRED.}\label{fig:m3:agriculturals-and-softs:cot}
\end{omfigure}

Dans les données hebdomadaires du maïs de janvier 2016 à septembre 2026, la \omterm{def:m3:agriculturals-and-softs:cot}{gestion pour compte de tiers} était nette acheteuse à un niveau record de 24.4~\%
des positions ouvertes le 22 mars 2022 et vendeuse nette de 22.8~\% le 9 juillet 2024~; elle a été vendeuse nette dans 300 des 559
semaines, et le 15 septembre 2026 elle était nette acheteuse de 22.5~\%, soit 2.26 écarts types au-dessus de sa moyenne des
trois années précédentes. Les traders lisent le rapport comme une mesure de la concentration~: un marché où les fonds sont
déjà très acheteurs~: il reste moins d'acheteurs en réserve si l'actualité déçoit.

\section{Softs et crush}

\begin{definition}[Softs]\label{def:m3:agriculturals-and-softs:softs}
Les \emph{softs}\index{softs} sont des matières premières agricoles cultivées plutôt qu'extraites, et qui ne sont
ni des céréales ni du bétail~: café, cacao, sucre, coton et jus d'orange.
\end{definition}

Les \omterm{def:m3:agriculturals-and-softs:softs}{softs} sont cultivés dans peu de pays, souvent par de petits exploitants, et leurs chocs d'offre sont concentrés. Le
Selon le FMI, le prix mensuel du cacao s'élevait en moyenne à 2\,369~\$ la tonne en 2022~; il a ensuite plus que quadruplé, pour atteindre une moyenne mensuelle de
10\,710~\$ en janvier 2025, et valait 5\,619~\$ en juillet 2026
(\cref{fig:m3:agriculturals-and-softs:cocoa}). L'International Cocoa Organization a indiqué que les \omterm{def:m3:agriculturals-and-softs:cropyear}{campagnes} 2022/23 et 2023/24
se sont achevées en déficit, la seconde à la suite d'aléas climatiques, de maladies et de ravageurs ayant réduit les récoltes des principaux producteurs d'Afrique de l'Ouest.

\begin{omfigure}
\begin{tikzpicture}
\begin{axis}[width=11cm,height=4.4cm,
  ylabel={milliers de \$ par tonne},
  tick label style={font=\small},label style={font=\small},
  xmin=2000,xmax=2026.75,ymin=0,ymax=12,ytick={5,10},grid=major,grid style={gray!25},
  xtick={2000,2005,2010,2015,2020,2025},xticklabel style={/pgf/number format/1000 sep={}},
  table/col sep=comma]
\addplot[omProp,very thick] table[x=t,y=cocoa]{figdata/markets-3/09-agriculturals-and-softs/cocoa.csv};
\end{axis}
\end{tikzpicture}

\omcaption{Cacao, prix mensuel moyen, de janvier 2000 à juillet 2026. Données~: FMI, Primary Commodity Prices
(série PCOCOUSDM via FRED).}\label{fig:m3:agriculturals-and-softs:cocoa}
\end{omfigure}

La transformation relie les cultures à leurs produits. Un broyeur de soja transforme les graines en tourteau (alimentation animale) et en huile
(alimentation humaine et carburant).

\begin{definition}[crush spread]\label{def:m3:agriculturals-and-softs:crush}
Le \emph{crush spread}\index{crush spread} est la valeur du tourteau et de l'huile issus d'un boisseau de graines de soja
diminuée du prix du boisseau. Aux rendements conventionnels de la bourse, soit 44 livres de tourteau et 11 livres
d'huile par boisseau de 60 livres, elle vaut, en dollars par boisseau,
$0.022 \times P_{\mathrm{meal}} + 0.11 \times P_{\mathrm{oil}} - P_{\mathrm{beans}}$, avec le tourteau exprimé en dollars
par short ton et l'huile en cents par livre.
\end{definition}

La marge de broyage est, pour un transformateur de soja, ce que le crack est pour un raffineur
(\cref{ch:m3:refined-products-and-cracks})~: sa marge, couverte en achetant du soja et en vendant du tourteau et de l'huile,
dans le ratio de dix contrats de soja pour onze de tourteau et neuf d'huile.

\section{Tutoriel~: positionnement et limites}\label{tut:m3:agriculturals-and-softs:cot}

\begin{tutorial}
\textbf{Objectif.} Transformer les fichiers COT en une mesure de positionnement, et calculer les limites et le blocage
d'un jour de rapport. \textbf{État final~:}
\cref{fig:m3:agriculturals-and-softs:cot,fig:m3:agriculturals-and-softs:limits} et la marge du
problème du week-end.
\begin{tutsteps}
\item \textbf{Positionnement.} Positions nettes, leur part des positions ouvertes et un z-score glissant.
\omcode{code/firm/cot/firm_cot.py}{12}{32}{Positions nettes et leur z-score glissant.}\label{lst:m3:agriculturals-and-softs:cot}
\item \textbf{Limites.} La règle de limite variable et l'évolution d'un \omterm{def:m3:agriculturals-and-softs:limit}{marché bloqué en limite}.
\omcode{code/firm/cot/firm_cot.py}{35}{63}{La règle de limite et une trajectoire de cours bloquée en limite.}\label{lst:m3:agriculturals-and-softs:limits}
\item \textbf{Exécuter} \texttt{m3\_ags.mm\_stats()}, \texttt{report\_day()} et \texttt{fig\_ags.py}.
\end{tutsteps}
\textbf{Pour aller plus loin.} Calculer la même mesure pour le soja et le blé, et vérifier si les positions des fonds
sur les trois marchés évoluent ensemble~; tester si un positionnement extrême prédit le rendement du mois suivant,
en corrigeant du chevauchement des échantillons hebdomadaires.
\end{tutorial}

\section{Construction~: l'analyseur des positions et l'état des limites}\label{bld:m3:agriculturals-and-softs:cot}

\begin{build}
\bfield{Rôle} Le firm miniature négocie les céréales autour des rapports~: il a besoin des positions hebdomadaires
par catégorie, des limites quotidiennes courantes de chaque contrat et de l'état d'un \omterm{def:m3:agriculturals-and-softs:limit}{marché bloqué en limite} pour
son système de risque et son routeur d'ordres, qui ne doit pas envoyer d'ordres qui ne peuvent pas être exécutés.
\bfield{Interface} \texttt{net(row, category)}, \texttt{net\_share}, \texttt{zscore(values, window)}~;
\texttt{variable\_limit(average\_price)}~; \texttt{limit\_path(start, fair, limit, expanded)}~;
\texttt{crush\_margin(beans, meal, oil\_cents)}.
\bfield{Règles} Les positions sont des entiers en contrats~; les z-scores n'utilisent que des valeurs passées~; la \omterm{def:m3:agriculturals-and-softs:limit}{limite
élargie} ne s'applique que le lendemain d'une clôture en limite~; les limites en dollars par boisseau, au cent près.
\bfield{Tests d'acceptation} \texttt{code/firm/cot/tests/}~: positions nettes~; la fenêtre du z-score~; la règle de
limite à trois niveaux de cours, y compris le floor~; un blocage de deux jours~; le crush du tableau.
\bfield{Pour aller plus loin} Transmettre les états de limite à \texttt{firm.calendar} (Livre 1, chapitre 22)~; lire
directement les fichiers COT dans le format d'archive de la Commission.
\end{build}

\begin{omsources}
\item US Commodity Futures Trading Commission, \emph{Commitments of Traders}~: à propos des rapports, du calendrier
de publication et des fichiers historiques désagrégés 2016--2026.
\item USDA, \emph{WASDE}~; National Agricultural Statistics Service, \emph{Acreage}, 30 juin 2026~; documentation ERS
des céréales fourragères (années de commercialisation).
\item CME Group, CBOT rulebook chapter 10 (corn) et notices de révision des limites de variation de cours des céréales~;
guide de référence du crush du soja.
\item International Monetary Fund, Primary Commodity Prices (cocoa, maize) via FRED.
\item International Cocoa Organization, \emph{Quarterly Bulletin of Cocoa Statistics}, novembre 2024 (résumé).
\item CME Group, price limits page for agricultural futures, trade date 8 septembre 2026 (copie de l'Internet Archive).
\end{omsources}

\section{Exercices}

\begin{exercise}[$\star$]\label{exo:m3:agriculturals-and-softs:1}
Les stocks de fin de \omterm{def:m3:agriculturals-and-softs:cropyear}{campagne} s'élèvent à 1.2 milliard de boisseaux et l'utilisation totale à 14.8 milliards. Calculer le \omterm{def:m3:agriculturals-and-softs:wasde}{ratio stocks/utilisation}.
\end{exercise}

\begin{exercise}[$\star$]\label{exo:m3:agriculturals-and-softs:2}
Donner les limites initiale et élargie pour du maïs dont le prix moyen est de 4.00~\$ et 6.00~\$ le boisseau.
\end{exercise}

\begin{exercise}[$\star$]\label{exo:m3:agriculturals-and-softs:3}
Pourquoi le maïs de juillet est-il un contrat de l'\omterm{def:m3:agriculturals-and-softs:cropyear}{ancienne récolte} et celui de décembre un contrat de la \omterm{def:m3:agriculturals-and-softs:cropyear}{nouvelle récolte}~?
\end{exercise}

\begin{exercise}[$\star\star$]\label{exo:m3:agriculturals-and-softs:4}
Le tourteau vaut 300~\$ le short ton, l'huile 50 cents la livre, le soja 10.00~\$ le boisseau. Calculer le crush du tableau.
\end{exercise}

\begin{exercise}[$\star\star$]\label{exo:m3:agriculturals-and-softs:5}
Un rapport fait baisser la juste valeur du maïs de 0.60~\$ avec des limites de 30 et 45 cents. Combien de jours une position longue
est-elle bloquée, et que perd-elle par contrat avant de pouvoir vendre~?
\end{exercise}

\begin{exercise}[$\star\star$]\label{exo:m3:agriculturals-and-softs:6}
Les fonds détiennent une position nette longue record juste avant un rapport décevant. Pourquoi la baisse pourrait-elle être
plus forte que ce que le seul rapport justifie~?
\end{exercise}

\begin{exercise}[$\star\star\star$]\label{exo:m3:agriculturals-and-softs:7}
\emph{Codage.} Avec \texttt{mm\_series}, donner la plus grande position nette longue et la plus grande position nette courte
de la \omterm{def:m3:agriculturals-and-softs:cot}{gestion pour compte de tiers} sur le maïs depuis 2016, leurs dates, et le nombre de semaines en position nette courte.
\end{exercise}

\begin{exercise}[$\star\star\star$]\label{exo:m3:agriculturals-and-softs:8}
\emph{Repérer l'erreur.} « La position nette des fonds et le cours évoluent ensemble, donc les fonds provoquent les mouvements
de cours~; nous allons trader devant eux. »
\end{exercise}

\section{Problème~: jour de rapport}

\begin{problem}[{Problème du week-end --- une position longue bloquée par un rapport}]\label{pb:m3:agriculturals-and-softs:1}
Un trader détient 50 contrats à terme longs sur le maïs à 4.50~\$ le boisseau. Un rapport ramène le prix de juste valeur à 3.70~\$. Les limites
suivent la règle de l'encadré daté pour un prix moyen de 4.50~\$.

\medskip\noindent\textbf{Partie I --- Les limites.}
\begin{enumerate}
\item Donner les limites initiale et élargie.
\item Quelle est la variation de valeur du contrat par boisseau~?
\item Où le maïs se règle-t-il le premier jour~?
\item Et le second~?
\item Quel jour le trader peut-il vendre, et à quel prix~?
\end{enumerate}

\medskip\noindent\textbf{Partie II --- L'argent.}
\begin{enumerate}[resume]
\item Quelle marge de variation le trader verse-t-il le jour 1~?
\item Et le jour 2~?
\item Quelle marge doit-il détenir pour survivre aux deux jours bloqués~?
\item Quelle est sa perte totale lorsqu'il vend le jour 3~?
\item Qu'aurait fait un ordre stop à 4.30~\$~?
\end{enumerate}

\medskip\noindent\textbf{Partie III --- Les couvertures.}
\begin{enumerate}[resume]
\item Pouvait-il se couvrir avec des options pendant le blocage~? À quel prix~?
\item Pouvait-il vendre un contrat apparenté (une autre échéance, une autre bourse)~?
\item Comment un producteur qui couvre sa récolte vit-il les mêmes journées~?
\item Pourquoi les bourses maintiennent-elles des limites malgré cela~?
\item Comment dimensionneriez-vous une position avant un rapport, compte tenu des limites~?
\end{enumerate}

\medskip\noindent\textbf{Partie IV --- Jugement.}
\begin{enumerate}[resume]
\item Que vous apprend le rapport COT sur le risque d'une telle journée de rapport~?
\item Pourquoi une limite n'est-elle pas un plafond de perte~?
\item Comment un blocage modifie-t-il la valeur de la liquidité~?
\item Énoncer le \emph{résultat nommé}~: la marge nécessaire pour survivre au blocage, par contrat et pour la position.
\item En une phrase~: qu'est-ce qu'une limite quotidienne~?
\end{enumerate}
\end{problem}

\section{Questions d'entretien}

\begin{interviewq}[$\star$ \iqroles{trader}]\label{iq:m3:agriculturals-and-softs:1}
Quel est le \omterm{def:m3:agriculturals-and-softs:wasde}{ratio stocks/utilisation} et pourquoi le cours en dépend-il de manière non linéaire~?
\end{interviewq}

\begin{interviewq}[$\star$ \iqroles{trader, chercheur}]\label{iq:m3:agriculturals-and-softs:2}
Qu'est-ce qu'un \omterm{def:m3:agriculturals-and-softs:limit}{marché bloqué en limite}, et comment gère-t-on le risque qu'il représente~?
\end{interviewq}

\begin{interviewq}[$\star\star$ \iqroles{chercheur}]\label{iq:m3:agriculturals-and-softs:3}
Comment construiriez-vous un signal à partir du rapport COT sans biais d'anticipation~?
\end{interviewq}

\begin{interviewq}[$\star\star$ \iqroles{trader}]\label{iq:m3:agriculturals-and-softs:4}
Expliquez le crush de soja et la manière dont un transformateur le couvre.
\end{interviewq}

\begin{interviewq}[$\star\star$ \iqroles{risque}]\label{iq:m3:agriculturals-and-softs:5}
Comment calculez-vous une mesure de type VaR pour une position sur un marché soumis à des limites quotidiennes~?
\end{interviewq}

\begin{interviewq}[$\star\star\star$ \iqroles{développeur, trader}]\label{iq:m3:agriculturals-and-softs:6}
Concevez le système de scénario pour les journées de rapport~: quelles données, quels contrôles, quelles actions automatiques à midi~?
\end{interviewq}
